\chapter{指数分布}
    \section{定義}
        $\mu>0$とする。
        確率密度関数$f(x)$が次式となる分布を、尺度母数$\mu$の指数分布という。
        \[
            f(x) =
            \begin{cases}
                \frac{1}{\mu}e^\frac{-x}{\mu} & (x\geq 0) \\
                0 & (x<0)
            \end{cases}
        \]
        確率変数$X$がこの分布に従うことを$X \sim \ExpDist{\mu}$と表すことにする。
    \section{解釈}
        単位時間当たりの生起回数の期待値が$\lambda$であるような事象が時間$x>0$の間に生起する確率の累積分布。
    \section{無記憶性からの導出}
        \begin{proof}
            \quad\par
            「観測開始から時間$s$の間イベントが起こらなかったという事象の下、それから時間$\Delta t$以内にイベントが起こる条件付き確率」が「観測開始($t=0$)から時間$\Delta t$以内にイベントが起こる確率」に等しいという条件より次式が成り立つ。
            \begin{align*}
                \frac{\integrate{s}{s+\Delta t}{f(x)}{}{x}}{\integrate{s}{\infty}{f(x)}{}{x}} &= \integrate{0}{\Delta t}{f(x)}{}{x} \\
                \therefore\; \integrate{s}{s+\Delta t}{f(x)}{}{x} &= \integrate{0}{\Delta t}{f(x)}{}{x} \integrate{s}{\infty}{f(x)}{}{x} \\
                \therefore\; f(s+\Delta t) - f(s) &= -f(s)\integrate{0}{\Delta t}{f(x)}{}{x} \quad (\text{両辺を}s\text{で微分した}) \\
                \therefore\; f'(s) &= -f(0)f(s) \quad (\text{両辺を}\Delta t\text{で割った後}\Delta t \to +0\text{とした})
            \end{align*}
            この微分方程式を解いて規格化すると$f(x) = f(0)\exp(-f(0)x)\;(f(0)>0)$となる。$\lambda = f(0)$とすれば$f(x) = \lambda e^{-\lambda x}$となる。
        \end{proof}
    \section{特性関数}
        \label{指数分布の特性関数}
        尺度母数$\mu$の指数分布の特性関数は次式である。
        \begin{align*}
            \E{}{e^{itX}} &= \integrate{0}{\infty}{\frac{1}{\mu}e^\frac{-x}{\mu}e^{itx}}{}{x} = \frac{1}{\mu}\integrate{0}{\infty}{e^{\left(\frac{-1}{\mu} + it\right)x}}{}{x} = \frac{1}{\mu}\frac{1}{\frac{-1}{\mu} + it}\left[e^{\left(\frac{-1}{\mu} + it\right)x}\right]_0^\infty = \frac{1}{1-i\mu t}
        \end{align*}